\scanpage{68}
\mspage{61}
% archive sheet number: 67
De plus, l'hom. canonique (pour $s\in\mathrm{Ob}\,I$, \inserted{$X'\in\mathrm{Ob}\,\mathrm{HO}_W(I')$})
\[
u_*^W(X')(s)\simeq p_{s*}\bigl(X'|_{\deleted{\unsure{s\backslash}}I'_{/s}}\bigr)
\overset{\text{def}}{=}
\underset{\deleted{\unsure{s\backslash}}I'_{/s}}{\underset{\longleftarrow}{\operatorname{\unsure{Holim|Rolim}}}}_W\bigl(X'|_{\deleted{\unsure{s\backslash}}I'_{/s}}\bigr)
\]
(où $p_s:\deleted{\unsure{s\backslash}}I'_{/s}\longrightarrow e$) est un iso, si on suppose que les $I'_{/s}$ sont admissibles \deleted{\unsure{(et $u$ admissible)}} \inserted{(\uline{strictement} admissible)}.

\uline{Corollaire} Sur le domaine formé des ens. ordonnés finis, \inserted{le prédérivateur} $\mathrm{HO}_W$ satisfait (en plus des axiomes $\mathrm{Der}\,1$, $\mathrm{Der}\,2$) aux axiomes $\mathrm{Der}\,4$ et $\mathrm{Der}\,5$. (existence des images directes $u_*$, adj. à droite des foncteurs $u^*$, et calcul \unsure{standard|standart} des fibres d'une image directe).

Démonstration du th.-scholie.

1°) a) est immédiat. Pour b) on procède par récurrence sur le cardinal $n$ de l'ens. ordonné fini $I$, les cas $n\in\{0,1\}$ étant triviaux. Supposons donc $n\geq2$ et le th. prouvé pour les $n'<n$. Je distingue deux cas : si $I$ admet un plus grand élém\textsuperscript{t} \deleted{\unsure{$b$}}\inserted{$a$}. Alors, désignant par $\alpha:e\unsure{-|\to}I$ le foncteur de valeur $a$, \inserted{par $p:I\unsure{-|\to}e$ la projection} on sait que
\[
p_*^{\deleted{\unsure{W|*}}}\simeq\alpha^{\unsure{*|\alpha}}\quad\text{i.e.}\quad\varprojlim_I X_i\simeq X_a,
\]
d'où le fait que sur $M(I)$ tout entier
